\begin{afterword}
\summaryhead

The author dates the end of the earlier folly to the moment of first grasping, in full generality, the notion of an optimization process. In 2002, writing ``Levels of Organization in General Intelligence'', the author saw natural selection and human intelligence as opposites, and still took the human mind as the starting point for AI. When Emil Gilliam asked whether the design hewed too close to the human one, the reply was that a completely alien design would be too hard to build.

What broke this was an unfinished science-fiction story with a non-cognitive optimizing effect. With three cases, the author generalized: an optimization process squeezes the future into a narrow region of the possible. This view goes beyond defining intelligence as ``achieving goals'', a phrase that invites us to imagine an agent that wants. Seeing an optimizer as a physical process that steers the future wherever its utility function points, the author saw that the 1997 design would have wiped out humanity, and concluded: ``I've been stupid.''

\responsehead

This is memoir about an insight, and the account of the author's own writings holds up where it could be checked. ``Levels of Organization'' was posted in April 2002 and published in 2007, as the post says. A footnote calls the comparison of evolution with intelligence a ``loose analogy'' rather than a ``special case'', which is the dichotomy the post describes. Legg's collection has the 71 definitions the post counts, and most of those by AI researchers do speak of goals or problems. The date of the awakening agrees with ``The Magnitude of His Own Folly'', five days later, which puts it in late 2002.

The one place where an old text adds to the story is the paperclip case. The post says that in 2001 its author did not know whether a superintelligence could pursue something pointless. \textsc{Creating Friendly AI} (2001) already describes such a case, an AI whose only supergoal is the Riemann Hypothesis, as one way to get a catastrophe. That does not settle what the memoir's author knew, but a reader should know the scenario was on the page.

The post presents the idea as a personal discovery, ``for me at least'', and it keeps that scope. The idea itself has predecessors worth knowing. Cybernetics gave ``goal'' a physical reading in 1943, as a final state of correlation between a system and its surroundings, which is close to the post's squeezing of the future; so the post's contrast is between a mentalistic and a physical reading of the word ``goal''. Legg's own collection contains a definition close to the post's, Lenat and Feigenbaum's ``power to rapidly find an adequate solution'' in ``an immense search space.'' And the lumping of thought with natural selection that the post dislikes is an established view, going back at least to Donald Campbell in 1960, and the post takes a side against it.

In short: a memoir whose account of its author's own papers checks out, and whose central idea has predecessors in cybernetics and in the collection it cites.
\end{afterword}
